\documentclass[conference]{IEEEtran}

\usepackage{amsmath}
\usepackage{booktabs}
\usepackage{graphicx}
\usepackage{hyperref}
\usepackage{xcolor}

\newcommand{\pending}{\textit{Pending controlled evaluation}}

\title{VeriSlip: Multi-Evidence Forensics for Bank-Transfer Slip Triage}

% MERCon 2026 requires an anonymous initial submission. Restore verified author
% names, affiliations, and acknowledgements only for a venue/stage that permits it.
\author{\IEEEauthorblockN{Anonymous Authors}}

\begin{document}
\maketitle

\begin{abstract}
Images of bank-transfer slips are frequently used as provisional payment proof,
yet small edits to amounts, beneficiaries, references, or status text can change
their meaning. Natural-image forensic assumptions transfer imperfectly to these
structured, text-heavy documents, while benign screenshotting, recompression,
and recapture can resemble manipulation. This paper presents VeriSlip, an
open-source decision-support architecture that combines structural and semantic
checks, classical compression and copy--move cues, local residual consistency,
and an optional dual-stream inference component. We formalize the observation
and attacker model, describe the implementation-aligned evidence fusion, and
specify a reproducible evaluation protocol for Sri Lankan mobile-payment slips.
Empirical values are deliberately omitted until the protocol is executed on a
versioned, consented dataset; therefore this artifact is a conference draft and
not an accuracy claim.
\end{abstract}

\begin{IEEEkeywords}
document forensics, receipt tampering, image manipulation, OCR, financial fraud
\end{IEEEkeywords}

\section{Introduction}
Merchants increasingly receive screenshots or photographs of transfer receipts
before releasing goods. A convincing false proof may require changing only a few
glyphs. The task is not equivalent to generic natural-image manipulation:
receipts have uniform backgrounds, repeated typography, rigid layouts, and
semantically constrained fields. At the same time, legitimate messaging and
capture pipelines add compression, scaling, moir\'e, or print--scan artifacts.

VeriSlip addresses image-level triage rather than transaction settlement. Its
central design premise is that no single artifact proves fraud. Instead, weak and
heterogeneous evidence is retained, localized where possible, and fused into an
actionable risk category. The intended contributions, subject to the evaluation
specified in Section~\ref{sec:evaluation}, are:
\begin{itemize}
  \item a financial-document threat model separating malicious edits from benign
  acquisition transforms;
  \item an implementation-aligned, multi-evidence architecture spanning document
  semantics and image forensics; and
  \item a reproducible protocol for document-, region-, calibration-, and
  generalization-level evaluation.
\end{itemize}

This draft makes no claim of official bank-network access, universal forgery
detection, or measured superiority over prior work.

\section{Related Work}
Passive image forensics studies traces left by acquisition and manipulation
pipelines~\cite{farid2009survey,verdoliva2020media}. JPEG history can be localized
through block-level artifacts~\cite{bianchi2012jpeg}; copy--move detection has
used invariant keypoints~\cite{amerini2011sift} and learned source/target
localization~\cite{wu2018busternet}. Sensor-pattern noise supports source-camera
identification when a reference fingerprint and suitable images exist
~\cite{lukas2006prnu}, assumptions that often fail for screenshots.

Document images add text and layout constraints. DocTamper targets tampered text
under document-specific compression conditions~\cite{qu2023doctamper}. VeriSlip
differs in scope: it combines deterministic financial semantics, classical image
cues, local residuals, and an optional learned component. Error-level analysis
(ELA) is used only as a recompression-difference feature and visualization; it is
not treated as a calibrated standalone detector.

\section{Problem Formulation and Threat Model}
Let a genuine slip be $G\in\{0,\ldots,255\}^{H\times W\times C}$. A manipulation
$M_\theta$ changes a material region $R$, producing $X=M_\theta(G;R)$. Genuine or
manipulated images then pass through a benign acquisition channel $T_\phi$:
\begin{equation}
  Y=T_\phi(X)+\epsilon,
\end{equation}
where $\epsilon$ represents sensor, display, compression, and OCR noise. For a
genuine observation, $X=G$ and $R=\varnothing$.

The adversary may replace amount, account, beneficiary, reference, date, or
status fields; copy, splice, erase, or inpaint regions; and resize, recompress,
screen-recapture, or print--scan the result. Novice through expert attackers are
considered, including adaptive probing. Genuine replay is a distinct risk: an
unaltered old receipt may be pixel-authentic yet irrelevant to the current order.

The verifier emits evidence $e$, candidate regions $\widehat R$, and bounded
risk $s=f(e)\in[0,1]$. Its verdict is image-forensic decision support, not proof
of payment finality. High-quality reconstruction, unknown templates, OCR error,
and severe recapture remain residual risks.

\section{Methodology}
\label{sec:method}

\subsection{Ingestion and normalization}
Public ingestion paths content-validate and fully decode bounded raster images,
reject unsupported formats, limit dimensions, detach decoded pixels, and remove
metadata during canonicalization. The analysis image is RGB with longest
dimension at most 1400 pixels.

\begin{figure*}[t]
\centering
\setlength{\tabcolsep}{4pt}
\renewcommand{\arraystretch}{1.35}
\begin{tabular}{c c c c c}
\fbox{\shortstack{Sanitized\\RGB input}} & $\rightarrow$ &
\fbox{\shortstack{L1/1.5\\structure + semantics}} & $\searrow$ & \\
&& \fbox{\shortstack{L2/2.5\\ELA, DCT, copy--move}} & $\rightarrow$ &
\fbox{\shortstack{weighted evidence\\+ gated maxima}} $\rightarrow$
\fbox{\shortstack{risk, verdict,\\regions, findings}}\\
&& \fbox{\shortstack{L3\\residual consistency}} & $\nearrow$ & \\
&& \fbox{\shortstack{L4 + adjuncts\\optional learned path}} & $\nearrow$ &
\end{tabular}
\caption{Executable evidence flow. The components run on the normalized image;
fusion preserves strong corroborated evidence rather than averaging it away.}
\label{fig:architecture}
\end{figure*}

\subsection{Structural and semantic evidence}
Layer 1 evaluates decoded metadata, aspect ratio, bank-color similarity, and an
optional reference format. With binary conditions $m,a,c,r$, its implemented
score is
\begin{equation}
s_1=\min(1,0.45m+0.15a+0.20c+0.40r).
\end{equation}
Local OCR supplies uncertain tokens and boxes to semantic rules for currency
grouping, mandatory fields, temporal order, and
$|\mathrm{total}-(\mathrm{amount}+\mathrm{fee})|>0.05$. The semantic score is
the maximum confidence assigned by a triggered rule.

\subsection{Classical and residual evidence}
Layer 2 re-encodes the image in memory as JPEG and computes
$D=|I-J_Q|$ plus $\mathrm{Var}(\mathrm{gray}(D))$. It analyzes 8-by-8 DCT AC
coefficient histogram periodicity, JPEG-grid phase, ORB and dense block-DCT
copy--move evidence, occlusion candidates, screen recapture, and thermal-fade
characteristics. These are heuristic cues with benign confounders.

Layer 3 forms $R=|G-\mathrm{median}_{3\times3}(G)|$ and calculates variance in
32-by-32 flat-background blocks. A block is a candidate when
\begin{equation}
z_b=\frac{v_b-\mu_{flat}}{\sigma_{flat}+10^{-4}}>3.5.
\end{equation}
Divider-like one-dimensional clusters are removed. The layer also summarizes
fixed spatial-rich-model residual kernels. This is local residual analysis, not
PRNU camera attribution.

\subsection{Optional inference and fusion}
The Layer 4 code defines RGB and forensic streams. The forensic tensor contains
ELA difference, median residual, and Sobel gradient channels. A compatible
checkpoint may produce classification and localization outputs. When no
checkpoint is loaded, the initialized-network path must not be interpreted as a
validated model; when PyTorch is unavailable, a statistical ELA/noise fallback
is used.

Default linear evidence is
\begin{equation}
L=0.18s_1+0.225s_2+0.225s_3+0.27s_4+0.10s_v,
\end{equation}
with a conditional recapture term. The final risk is the maximum of $L$ and
corroboration-gated peak signals. Candidate regions are merged by greedy
non-maximum suppression at intersection-over-union 0.3. This score is a bounded
triage heuristic, not a posterior probability.

\section{Evaluation Protocol}
\label{sec:evaluation}
The evaluation must use a versioned, authorized, privacy-preserving dataset. The
split unit should prevent bank-template, source-slip, merchant, and device
leakage. Genuine controls must include resize, repeated JPEG, messaging-app
transport, screen capture, camera recapture, and print--scan. Manipulations must
cover field replacement, erasure/inpainting, copy--move, splicing, and replay.

\begin{figure}[t]
\centering
\fbox{\parbox{0.92\columnwidth}{\centering
Versioned manifest $\rightarrow$ leakage-aware split $\rightarrow$ locked
configuration $\rightarrow$ document and region metrics $\rightarrow$ bootstrap
confidence intervals $\rightarrow$ error analysis}}
\caption{Required reproducible evaluation sequence.}
\label{fig:evaluation}
\end{figure}

Report precision, recall, F1, ROC-AUC and PR-AUC at document level; IoU and
pixel-level measures for localized edits; expected calibration error and
abstention coverage; latency and memory; and per-bank, language, device,
acquisition, and attack breakdowns. Thresholds are chosen on validation data and
locked before test evaluation. Bootstrap confidence intervals and paired
comparisons should accompany point estimates.

\begin{table}[t]
\caption{Evidence status in this draft (no values are inferred from demos).}
\label{tab:results}
\centering
\begin{tabular}{lp{0.44\columnwidth}}
\toprule
Artifact & Status \\
\midrule
Dataset manifest and consent record & \pending \\
Locked train/validation/test split & \pending \\
Document and localization metrics & \pending \\
Layer ablations and baselines & \pending \\
Confidence intervals and error audit & \pending \\
\bottomrule
\end{tabular}
\end{table}

The repository contains deterministic demonstration values and synthetic test
fixtures. They validate software paths but are not experimental observations and
are excluded from Table~\ref{tab:results}.

\section{Discussion}
The architecture makes heterogeneous evidence inspectable, but fixed rules and
max fusion can be miscalibrated across banks or acquisition channels. OCR and
template checks may fail on multilingual or redesigned interfaces. Compression
and recapture can either create false positives or erase manipulation traces.
The optional learned component requires a documented checkpoint and independent
test set before supporting an accuracy claim. Operationally, an `AUTHENTIC`
image verdict must never substitute for checking account balance or an
authorized transaction record.

Privacy is part of the research design: collect only authorized slips, minimize
retention, redact research exports, separate identifiers from labels, restrict
access, and never place raw images or OCR values in logs or publications.

\section{Conclusion}
VeriSlip frames payment-slip screening as multi-evidence document forensics. Its
implementation combines semantic constraints with compression, copy--move,
residual, typography, recapture, and optional learned evidence. This draft
defines the system and a falsifiable evaluation plan. Completion of the
registered protocol, rather than repository demos, is required before any
performance or comparative conclusion.

\bibliographystyle{IEEEtran}
\bibliography{references}

\end{document}
